\begin{proposition}
\label{P:pushforward is local}
\begin{enumerate}
\item 设 $\mu_1$、$\mu_2$ 是支配余权，$\nu$ 是支配余权。
设 $(m_1,n_1,m_2,n_2)$ 是 $(\mu_1+\nu, \nu)$-可接受且 $(\mu_2+\nu,\nu)$-可接受的四元组。
则存在 Hecke 对应
\begin{equation}\label{deompf1}
\Sht_{\mu_1}^{\loc(m_1,n_1)}\xleftarrow{\overleftarrow{h}_{\mu_1}^{\loc(m_1,n_1)}} \Sht_{\mu_1\mid \mu_2}^{\nu,\loc(m_1,n_1)}\xrightarrow{\overrightarrow{h}_{\mu_2}^{\loc(m_2,n_2)}} \Sht_{\mu_2}^{\loc(m_{2},n_{2})},
\end{equation}
其中 $\Sht_{\mu_1|\mu_2}^{\nu,\loc(m_1,n_1)}$ 是代数叠，且与 $(m_1,n_1)$ 的选取无关。

此外，若 $(m'_1,n'_1,m'_2,n'_2)\geq (m_1,n_1,m_2,n_2)$ 是满足同样条件的另一个四元组（允许 $(m'_1,n'_1,m'_2,n'_2)=(\infty,\infty,\infty,\infty)$），
则下图交换，且左方块是笛卡儿的。
\begin{equation*}
\label{E:cycle construction diagram full}
\xymatrix@C=17pt{
\Sht_{\mu_1}^{\loc(m'_1,n'_1)}\ar_{\res^{m'_1,n'_1}_{m_1,n_1}}[d]
&\Sht_{\mu_1|\mu_2}^{\nu,\loc(m'_1,n'_1)}\ar_{\res}[d]\ar[l]\ar[r]& \Sht_{\mu_2}^{\loc(m'_2,n'_2)}\ar^{\res^{m'_2,n'_2}_{m_2,n_2}}[d]
\\
\Sht_{\mu_1}^{\loc(m_1,n_1)}&\Sht_{\mu_1|\mu_2}^{\nu,\loc(m_1,n_1)}\ar[l]\ar[r]& \Sht_{\mu_2}^{\loc(m_2,n_2)}.
}
\end{equation*}

\item 设 $\mu_1,\mu_2,\mu_3,\nu_1,\nu_2$ 是支配余权，$(m_1,n_1,m_2,n_2,m_3,n_3)$ 是六元组，使得每个四元组 $(m_i,n_i,m_{i+1},n_{i+1})$ 都是 $(\mu_i+\nu_i,\nu_i)$-可接受且 $(\mu_{i+1}+\nu_i,\nu_i)$-可接受的。则存在自然的完全紧合态射
\begin{equation*}
\label{E:composition restricted}
\on{Comp}_{\mu_1, \mu_2, \mu_3}^{\mathrm{loc}(m_1,n_1)}: \Sht_{\mu_1|\mu_2}^{\nu_1,\loc(m_1,n_1)} \times_{\Sht_{\mu_2}^{\loc(m_2,n_2)}} \Sht_{\mu_2|\mu_3}^{\nu_2,\loc(m_2,n_2)} \longto \Sht_{\mu_1|\mu_3}^{\nu_1+\nu_2,\loc (m_1,n_1)}.
\end{equation*}

此外，若 $(m'_1,n'_1,m'_2,n'_2,m'_3,n'_3)\geq (m_1,n_1,m_2,n_2,m_3,n_3)$ 是满足同样条件的另一个六元组（允许 $(m'_1,n'_1,m'_2,n'_2,m'_3,n'_3)=(\infty,\infty,\infty,\infty,\infty,\infty)$），则下图交换。
{\footnotesize
\begin{equation*}
\label{E:projection localization truncated}
\xymatrix@C=5pt{
&&\Sht_{\mu_1|\mu_3}^{\nu_1+\nu_2,\loc(m'_1,n'_1)}
\ar[ddd]\ar[dll]\ar[dr]
\\
\Sht^{\loc(m'_1,n'_1)}_{\mu_1} \ar[d]^{\res^{m'_1,n'_1}_{m_1,n_1}} &\ar[l]\ar[d] \Sht_{\mu_1\mid\mu_2}^{\nu_1,\loc(m'_1,n'_1)} \times_{\Sht_{\mu_2}^\loc(m'_2,n'_2)} \Sht_{\mu_2\mid\mu_3}^{\nu_2,\loc(m'_2,n'_2)} \ar[l] \ar[rr] \ar[ur]_-{\mathrm{Comp}^{\loc(m'_1,n'_1)}}&& \Sht^{\loc(m_3,n_3)}_{\mu_3} \ar[d]_{\res^{m'_3,n'_3}_{m_3,n_3}}
\\
\Sht_{\mu_1}^{\loc(m_1,n_1)}
&\ar[l]\Sht_{\mu_1|\mu_2}^{\nu_1,\loc(m_1,n_1)} \times_{\Sht_{\mu_2}^{\loc(m_2,n_2)}} \Sht_{\mu_2|\mu_3}^{\nu_2,\loc(m_2,n_2)} \ar[rd]^-{\mathrm{Comp}^{\loc(m_1,n_1)}} \ar[rr]
&&\Sht_{\mu_3}^{\loc(m_3,n_3)}
\\
&& \Sht_{\mu_1|\mu_3}^{\nu_1+\nu_2, \loc(m_1,n_1)}\ar[ull]\ar[ru].
}
\end{equation*}}
此外，中间的梯形是笛卡儿的。

\item 设 $\mu_1, \mu_2, \mu_3, \mu_4, \nu_1, \nu_2, \nu_3$ 是支配余权，$(m_1, n_1, m_2,n_2, m_3,n_3, m_4,n_4)$ 是非负整数八元组，使得 $(m_i,n_i,m_{i+1},n_{i+1})$ 是 $(\mu_i+\nu_i,\nu_i)$-可接受且 $(\mu_{i+1}+\nu_i,\nu_i)$-可接受的。则
\[\on{Comp}_{\mu_1,\mu_3, \mu_4}^{\mathrm{loc}(m_1,n_1)}\circ (\on{Comp}_{\mu_1,\mu_2, \mu_3}^{\mathrm{loc}(m_1,n_1)} \times \id)=\on{Comp}_{\mu_1,\mu_2, \mu_4}^{\mathrm{loc}(m_1,n_1)}\circ (\id \times \on{Comp}_{\mu_2,\mu_3, \mu_4}^{\mathrm{loc}(m_2,n_2)}).\]
\end{enumerate}
\end{proposition}

\begin{remark}
\label{R: grp theoretical Sht}
我们给出限制局部 $G$-Shtuka 模空间的群论解释。
给定 $(m,n)$-限制局部 $G$-Shtuka，把 $\calE_\rightone|_{D_{m}}$ 与 $\calE_\leftone|_{D_{n}}$ 都平凡化，则同构 $\psi$ 定义 $K_n$ 中的元素。
注意 $\Gr_\mmu^{(n)}$ 正是分类 $\Hk^{\loc(m)}_\mmu$ 上 $\calE_\rightone|_{D_{m}}$ 与 $\calE_\leftone|_{D_{n}}$ 两者的平凡化的模空间。于是
\[
\Sht_{\mmu}^{\loc(m,n)}\cong \frac{\Gr^{(n)}_{\mmu}\times^{K_n,\sigma} K_n}{\Ad K_m},
\]
其中右端表示 $\Gr_{\mmu}^{(n)}\times K_n$ 关于 $K_m\times K_n$ 的如下作用的商：
$$(k_m,k_n)\cdot (x,k)=(k_mxk_n^{-1},\sigma(k_n)kk_m^{-1}), \quad (k_m,k_n)\in K_m\times K_n, \ (x,k)\in \Gr_\mmu^{(n)}\times K_n.$$
特别地，$\Gr_{\mu}^{(n)}$ 的 Frobenius 自同构诱导同构
\[\sigma: \Sht_\mu^{\loc(m,n)}\cong \frac{\Gr_{\mu}^{(n)}}{\Ad_\sigma K_m}.\]
（在等特征情形，这些空间有典范的去完全化，且 $\Gr_{\mu}^{(n)}$ 的 Frobenius 自同态诱导 $\Sht_\mu^{\loc(m,n)}$ 与 $\frac{\Gr_{\mu}^{(n)}}{\Ad_\sigma K_m}$ 之间的泛同胚。）
因此，后文定义 $\Sht^\loc_{\bar k}$ 上反常层范畴时，既可以用 $\Sht_{\mmu}^{\loc(m,n)}$，也可以用 $\frac{\Gr^{(n)}_{\mmu}\times^{K_n,\sigma} K_n}{\Ad K_m}$。不过，后者（从群论上看）似乎更易描述，而前者才是与 Shimura 簇（或整体 Shtuka 的模空间）相联系所需要的。对 $\Sht_{\mu_1\mid \mu_2}^{\nu,\loc(m_1,n_1)}$ 也有类似的群论描述，细节参见 \cite[注 5.3.15]{XZ}。
\end{remark}

\subsection{范畴 $\on{P}^\Hk(\Sht^\loc_{\bar k})$}
有了上一小节的准备，就可以定义范畴 $\on{D}^\Hk(\Sht^\loc_{\bar k})$；由引理 \ref{L: fully faithful}，它可以看作 $LG$ 上 $\Ad_\sigma LG$-等变层范畴（的某个子范畴）的替代物。事实上，我们将主要关注子范畴 $\on{P}^{\Hk}(\Sht^\loc_{\bar k})\subset \on{D}^\Hk(\Sht^\loc_{\bar k})$，它将成为几何 Satake 的范畴迹的接收方。

首先，局部 Shtuka 模空间（基变换到 $\bar k$）上反常层的范畴 $\mathrm{P}(\Sht_{\bar k}^\loc)$ 容易定义。
设 $\mmu$ 是支配余权序列。对非负整数对 $(m',n')$ 与 $(m,n)$，若 $m \leq m'$、$n\leq n'$，且 $m'-n'$ 与 $m-n$ 都是 $\mmu$-大（$m'\neq \infty$），则层的平移 $!$-拉回函子
$$\Res^{m',n'}_{m,n}=(\res^{m',n'}_{m,n})^![((m'-n')-(m-n))\dim G]:\on{P}(\Sht^{\loc(m,n)}_{\mu})\to \on{P}(\Sht^{\loc(m',n')}_{\mu})$$
是反常精确的。由 \eqref{true-to-trun}，有函子的典范同构
\[\Res^{m',n}_{m,n}\circ\Res^{m',n'}_{m',n}\cong \Res^{m',n'}_{m,n}.\]
若 $m\geq 1$，则函子 $\Res^{m',n}_{m,n}$ 诱导范畴的等价。另一方面，$\res^{m',n'}_{m',n}$ 是 $\Ker\big( \Aut(\calE_\rightone|_{D_{n'}}) \to \Aut(\calE_\rightone|_{D_{n}}) \big)$ 下的主丛；若 $n\geq 1$，后者又是某个幂幺（非常值）概形的完全化，因此 $\Res^{m',n'}_{m',n}$ 全忠实。

于是可以把 $\Sht^\loc$ 上反常层的范畴定义为（直和意义下的）滤过极限
\begin{equation}
\label{E:definition of P(Sloc)}
\mathrm{P}(\Sht_{\bar k}^\loc): =\varinjlim_{(\mu, m,n)} \mathrm{P}(\Sht_\mu^{\loc(m, n)}),
\end{equation}
其中极限取遍三元组 $\{(\mu, m, n) \in \xcoch(T)/W \times \ZZ^2_{\geq 0} \mid m-n \mbox{ 是 }\mu\mbox{-大}\}$（带积偏序），
连接函子由全忠实函子的如下复合给出
\[
\mathrm{P}(\Sht_{\mu}^{\loc(m, n)})
\xrightarrow{\Res^{m',n'}_{m,n}}
\mathrm{P}(\Sht_{\mu}^{\loc(m', n')})
\xrightarrow{i_{\mu,\mu',*}}
\mathrm{P}(\Sht_{\mu'}^{\loc(m' n')}),
\]
其中 $i_{\mu, \mu'}: \Sht_{\mu}^{\loc(m', n')} \to \Sht_{\mu'}^{\loc(m', n')}$ 是自然闭嵌入。
注意由紧合或光滑基变换，这些连接函子满足自然的相容性条件，因此该极限确实有意义。
